\section{Leading-energy reduction of the central quadratic residue}
\label{sec:r13-leading-energy-residual}

This section follows PA and the protected background trace theorem
MT.  The space is the supplied original-field normal-three model,
the scalar presentation is the fixed finite continuous closed-cell
PL/PS class, and $P$ is the actual quadratic zero of PA1.  For an
arbitrary allowed degree-one background $S$, PA has proved
\[
                     \mathcal A_S(P)=c_S(P)\delta_y.
 \tag{LE1}
\]
The purpose here is to remove all higher coefficient jets from
this scalar remainder and to state its exact paired finite-
dimensional obstruction.  No trace of $S$ on a curve, no tangent
current of $S$, and no weak approximation of $S$ by classical
potential backgrounds will be assumed.

\subsection{Protected scalar first jets and all-face control}

Choose once the scalar ambient extensions supplied by SE, after
a common finite coefficient refinement.  Translate their central
logarithmic argument to zero.  Every cell whose closure does not
contain zero can be omitted after one shrinking.  The resulting
finite star consists of rational polyhedral cones.  If $h$ is
one of the continuous scalar extensions and $h_C$ is its smooth
extension on a closed cell $C$, define
\[
                 h^{[1]}(v)=d h_C(0)(v),\qquad v\in C.
 \tag{LE2}
\]
The constant $h(0)$ is omitted.  Differentiating the equality of
two scalar restrictions on a common face shows that LE2 agrees
there.  Thus $h^{[1]}$ is an actual continuous homogeneous PL
scalar in the old original moments.  Real coefficient slopes
are allowed, as in the original signed scalar class.  Its
piecewise differential has the prescribed tangential restriction
on every face.  The same construction applies to every $k_\nu$.

Set
\[
 Z_0=\tfrac12\sum_\nu c_\nu
   \left(d'_P h_\nu^{[1]}\wedge d''_P k_\nu^{[1]}
       +d'_P k_\nu^{[1]}\wedge d''_P h_\nu^{[1]}\right).
 \tag{LE3}
\]
This is a genuine regular finite-PS $(1,1)$ form in the
degree-one module, with constant coefficient tensors on the
cones and their consistent tangential face pullbacks.  It is
not a collection of independent coefficients on selected source
cells.  Choosing scalar extensions fixes the construction; only
their restrictions on realized directions enter LE3.  In
particular any two such extensions agreeing on the complete
ordinary image have the same directional tangent restrictions
on its finite central star.

Apply GP on a prescribed ordinary neighborhood where these
scalar identities hold, retaining the complete scalar tuple.
Write $F$ for the protected tuple, $b=t_F(y)$, and
$z=t_F-b$.  For a sufficiently small target ball, its entire
closed inverse image is compact in the original analytic
interior.  The whole central fiber, including points other
than $y$ and all lower moment ranks, is retained.  Put
\[
              \beta=\sum_{j=1}^{N}d'z_j\wedge d''z_j.
 \tag{LE4}
\]
Here and below a displayed target form denotes its original
moment pullback when used on the analytic space.

\begin{lemma}[Leading-energy error in the full regular order module]
\label{lem:r13-leading-energy-full-order}
There are constants $C$ and $r_0>0$, depending only on the
fixed finite scalar presentation and the protection, such that
on the complete inverse image of $|z|\le r<r_0$,
\[
                  -Cr\,\beta\ \le Z_P-Z_0\le Cr\,\beta.
 \tag{LE5}
\]
These are regular coefficient inequalities in the original
middle positive cone.  They pull back through every actual
refinement, including rank loss, and give the corresponding
original order inequalities after multiplication by every
allowed background.  No source point or charged coefficient
face has been removed.
\end{lemma}

\begin{proof}
For every fixed closed cell, Taylor's formula and its bounded
second derivatives give
\[
   |d h_C(v)-d h_C(0)|\le C_h|v|,
              \qquad |d h_C(v)|+|d h_C(0)|\le C_h.
 \tag{LE6}
\]
The finite maximum of the constants works for the whole star.
The old scalar coordinate vector is a fixed linear projection
of $z$, so its norm is at most a fixed multiple of $|z|$.
Subtracting the two products in the energy therefore gives
an operator norm bound $C|z|$ for the symmetric tensor on
each closed cell.  All tangential face restrictions are the
same restrictions of the scalar first jets used in LE2.
Thus the bound also holds on every lower face, even where
normal derivatives from different cells do not agree.

An operator norm bound relative to the Euclidean form in LE4
is precisely LE5.  Pullback by a linear tangent map preserves
both semipositive inequalities; it does not require an inverse
or a nonzero rank.  The original all-background regular-PS
order theorem therefore gives the asserted module inequalities
on all original tests.  The estimate is evaluated at every
ordinary point of the complete protected inverse image.  In
particular, the other points of its central fiber are not
discarded by using the distinguished point $y$.
\end{proof}

\subsection{Only the leading tensor contributes at the critical scale}

Let $\rho\in C_c^\infty(\mathbf R^N)$ be one fixed function
equal to one on a neighborhood of the closed unit ball and
supported in the ball of radius two.  For sufficiently small
$\delta>0$ put
\[
                  \chi_\delta=\rho(z/\delta).
 \tag{LE7}
\]
Protection makes this a compact original smooth scalar test,
without multiplying it by an additional source cutoff.  It is
one near the entire central fiber, in particular near $y$.

MT defines the positive Radon trace measure on the real target
by compact original pullback tests.  Denote it here by
$\mu_{S,F}$, with the harmless fixed normalization
$\mu_{S,F}(B)=\langle S,1_{z^{-1}(B)}\beta^2/2\rangle$.
MT proves that
\[
       \frac{\mu_{S,F}(B_r)}{r^2}\quad\hbox{is nondecreasing
       in }r,
       \qquad
       \nu_{S,F}:=\lim_{r\downarrow0}
                       \frac{\mu_{S,F}(B_r)}{r^2}<\infty.
 \tag{LE8}
\]
The convention for the factor $1/2$ only rescales this trace
and its density together.  MT uses the full compact protected
fiber, so separate closedness is not lost to a source cutoff.
The density in LE8 need not be zero.

\begin{theorem}[First-jet determination of the central action]
\label{thm:r13-leading-energy-residual}
For every allowed degree-one $S$ in the unchanged original
scope,
\[
 c_S(P)=-\lim_{\delta\downarrow0}
           \left\langle M_S(Z_0),D\chi_\delta\right\rangle.
 \tag{LE9}
\]
More precisely, for some fixed $C_P$ and every sufficiently
small $\delta$,
\[
 \left|c_S(P)+\langle M_S(Z_0),D\chi_\delta\rangle\right|
       \le C_P\delta^{-1}\mu_{S,F}(B_{2\delta})=O_S(\delta).
 \tag{LE10}
\]
The tensor $Z_0$ in this formula is the full paired energy
tensor in LE3.  Its two marginal families cannot be tested
or saturated separately.
\end{theorem}

\begin{proof}
The quadratic transgression gives
$\mathcal A_S(P)=-D M_S(Z_P)$: the other term
$\tfrac12D_\delta\sum c_\nu h_\nu k_\nu$ is already full
first-order exact and is killed after applying $D M_S$.
Pairing LE1 with LE7 therefore yields
\[
                    c_S(P)=-\langle M_S(Z_P),D\chi_\delta\rangle.
 \tag{LE11}
\]
This equality uses the current derivative sign of the retained
QA convention.  It holds for every sufficiently small $\delta$,
not just a subsequence.

The target Hessian of $\rho(z/\delta)$ has coefficient norm
at most $C_\rho\delta^{-2}$ and is supported in $B_{2\delta}$.
By LE5 the coefficients of $Z_P-Z_0$ are at most $C_P\delta$
there.  The matrix of their wedge product on
$\bigwedge^2\mathbf R^N$ is consequently bounded in norm by
$C\delta^{-1}$.  The form $\beta^2$ is strictly middle
positive on this exterior square.  It follows that
\[
 \pm\bigl((Z_P-Z_0)\wedge D\chi_\delta\bigr)
       \le C\delta^{-1}1_{B_{2\delta}}\beta^2.
 \tag{LE12}
\]
The bound is meant coefficientwise with the actual finite-PS
face values; it can equally be paired as a Borel Radon bound.
Every later tangent pullback preserves the inequality.
Multiplication by the original middle-positive $S$ and the
canonical regular-PS module action prove LE10.  This uses no
identification of middle, weak and strong positivity.
Finally LE8 implies $\mu_{S,F}(B_{2\delta})=O_S(\delta^2)$,
so the right side tends to zero.  All wall and central-fiber
contributions have been retained in this estimate.
\end{proof}

\begin{corollary}[Two precise sufficient central cancellations]
\label{cor:r13-leading-energy-cancellations}
Each of the following hypotheses implies $c_S(P)=0$ and
hence the full action identity on $U$ in PA:
\begin{enumerate}
\item the protected moment density $\nu_{S,F}$ is zero;
\item the leading energy $Z_0$ has an actual full-exact
degree-one identity on an ordinary neighborhood containing
the protected central fiber,
\[
                         Z_0=d'\beta_0+d''\zeta_0.
 \tag{LE13}
\]
\end{enumerate}
The second is a condition only on the finite homogeneous
leading tensor.  It does not require a full-exact identity
or a single controlled approximating sequence for the original
variable tensor $Z_P$.
\end{corollary}

\begin{proof}
The same matrix bound without the $O(\delta)$ error gives
$|\langle M_S(Z_0),D\chi_\delta\rangle|
\le C\delta^{-2}\mu_{S,F}(B_{2\delta})$.
In the first case its limit is zero.  In the second case QA's
first-derivative commutation and nilpotence give
$D M_S(Z_0)=0$, so its pairing is zero for every small
$\delta$.  Apply LE9.  PA already gives the punctured
vanishing required to conclude on the whole domain.
\end{proof}

\begin{corollary}[A concrete all-background first-jet range]
\label{cor:r13-linear-first-jet-quadratic-action}
Suppose every $h_\nu^{[1]}$ and $k_\nu^{[1]}$ is the
restriction of a single ambient linear form.  This holds, in
particular, when the closed-cell first derivatives of each
fixed scalar extension agree at the centre.  Then
\[
                         \mathcal A_S(P)=0
 \tag{LE13a}
\]
on the supplied ordinary model germ for every allowed original
background and every legal degree, on all original compact
smooth tests and actual refinements.  Higher scalar jets may
remain only finite PS, and no primitive for the full variable
energy $Z_P$ is required.  The first derivatives need agree
only at the centre: on punctured coefficient walls their
jumps may be nonzero, with size $O(|z|)$.  No $C^1$
regularity on a whole neighborhood is imposed.  The globally
smooth or globally $C^1$ subclass is not claimed as new;
the deduction here is from this pointwise central jet
condition and the PA--MT--LE composition.

More generally the same conclusion holds whenever the
\emph{combined} leading tensor $Z_0$, on every actual face,
is the restriction of one ambient constant symmetric tensor,
or differs from such a tensor by certified actual homogeneous
full-exact symbols.  The individual scalar first jets need
not be linear in this more general criterion.
\end{corollary}

\begin{proof}
Under the first hypothesis LE3 is a constant ambient symmetric
tensor $A$.  The ambient covector $b(z)=Az$ has
$e(b)=A$.  In the convention of QA this gives the concrete
full-exact identity
\[
 A=d'\!\left(\tfrac12\sum_j b_j(z)d''z_j\right)
        +d''\!\left(-\tfrac12\sum_j b_j(z)d'z_j\right).
\]
It is an identity on the whole ordinary moment domain and
survives every rank-losing tangent pullback.  Apply the second
case of the preceding corollary for degree-one backgrounds.
The lower-degree conclusion is already PA.  The proof is
unchanged when only the combined paired tensor is constant,
or when a finite sum of certified full-exact symbols is
subtracted from it first.
\end{proof}

\subsection{The finite paired residual functional}

Choose a finite-dimensional real vector space $\mathcal V$
of actual regular homogeneous degree-zero finite-PS $(1,1)$
forms in the protected coordinates, containing $Z_0$.  For
example take the span of the genuine paired product symbols
formed from the finite first-jet scalar list in LE2, together
with finitely many selected certified homogeneous exact symbols,
and then impose actual degree-one equality.  This construction
does not presume that an arbitrary face-compatible tensor
collection has an actual ambient presentation.  Fix a norm
on this finite coefficient space.  Let
$\mathcal E\subset\mathcal V$ be any linear subspace whose
members have certified actual identities of the form LE13 on
the complete protected neighborhood.  Ambient constant tensors
are examples.  A finite continuous piecewise-linear ambient
covector with its full all-face strain identity supplies another
certified member.  A collection of independent branch covectors
without joint compatibility does not.  The exact identities may
be given on an ordinary neighborhood of the entire compact
central fiber: properness then puts the support of every
sufficiently small LE7 cutoff in that neighborhood.  Indeed,
otherwise a sequence in its compact complement with target
values tending to zero would have a central-fiber limit in
that complement.

Define
\[
 \lambda_{\delta,S}(E)
       =-\langle M_S(E),D\chi_\delta\rangle,
                        \qquad E\in\mathcal V.
 \tag{LE14}
\]

\begin{proposition}[A full paired finite residual criterion]
\label{prop:r13-paired-finite-residual}
The family LE14 is bounded in $\mathcal V^*$ and vanishes
on $\mathcal E$.  Every cluster functional $\lambda_S$ satisfies
\[
 \begin{gathered}
 \lambda_S\in(\mathcal V/\mathcal E)^*,\qquad
                  \lambda_S([Z_0])=c_S(P),\\
 |\lambda_S([E])|\le C\nu_{S,F}\,
                 \operatorname{dist}(E,\mathcal E).
 \end{gathered}
 \tag{LE15}
\]
Consequently all cluster functionals have the same value on
the particular paired class $[Z_0]$, even if their values on
other classes differ.  Central cancellation is exactly their
annihilation of this class.  No limit current of $S$ is needed
to state or prove this criterion.
\end{proposition}

\begin{proof}
On the finite-dimensional space $\mathcal V$, the maximum
coefficient norm on the finite cone list is bounded by a
constant times the chosen norm.  The matrix estimate used
above gives
\[
        |\lambda_{\delta,S}(E)|
          \le C\|E\|\delta^{-2}\mu_{S,F}(B_{2\delta}).
 \tag{LE16}
\]
MT makes this family bounded, so finite-dimensional compactness
gives subsequential cluster functionals.  For $E\in\mathcal E$,
QA makes $D M_S(E)=0$ on the ordinary domain containing
the support of every sufficiently small cutoff, so LE14 is
zero.  It follows that all these functionals, and their limits,
factor through the quotient.  Apply LE16 to $E-E_0$ for
$E_0\in\mathcal E$, pass to the limit using LE8, and take
the infimum.  The fixed factor from $2\delta$ is absorbed
in $C$.  LE9 proves their common value on $Z_0$.
\end{proof}

\paragraph{Exact remaining gate and protected scope.}
The new conclusion is the unconditional $O(\delta)$ removal
of the higher scalar jets, after PA and MT.  It does not
show that the leading formal quadratic presentation is itself
an actual analytic zero; that stronger linearization assertion
is not used.  It does not identify the realized residual
functionals in LE15 with the full dual of a fixed old Killing
obstruction space.  The retained PR/QS examples already warn
against that identification: legal zero-preserving action data
do not read every nonzero fixed-old stress.  Nor do the available
curve caps prove that their realized paired rows span the
required class in every later actual refinement.

The all-background remaining statement is that the realized
functionals LE15 annihilate $[Z_0]$ for every actual zero in
PA1.  Positivity and separate closedness supply the critical
finite bound LE8, and actual full-exact symbols supply its
known annihilator.  They do not supply that final paired
annihilation.  In particular $O(\delta^2)$ is not replaced
by $o(\delta^2)$, and no restriction of the current order to
a selected curve or to a radial source family is performed.

\paragraph{Dependencies.}
The new analytic input beyond PA is
\texttt{joint13\_protected\_background\_trace\_and\_tangents.tex},
MT2--MT8, which derives the protected full-fiber test functional
and its radial mass bound from original tests and the real
supercurrent monotonicity theorem.  The scalar and regular-order
inputs are SE and the frozen R6 regular-PS module theorem.
The exactness implication is QA4/QA10.  All original fields,
norms, actual refinements, and absolute carriers are unchanged.
